\documentclass[../../main.tex]{subfiles}% 注意这里的文件路径不能用 ./main.tex ,否则用latexmk编译子文件会报错
\graphicspath{{\subfix{./image/}}{\subfix{../../image/}}} % 指定图片引用路径,后续可以直接使用图片文件名
% 如果图片存放在上级目录,则使用 ../../image/ 作为备用路径

% 例如:
% \begin{figure}[H]
% \centering
% \includegraphics[width=0.3\textwidth]{图.png}
% \caption{}
% \label{figure:图}
% \end{figure}
% 注意:上述\label{}一定要放在\caption{}之后,否则引用图片序号会只会显示??.

\begin{document}

\section{$L^p$空间的定义与不等式}

\begin{definition}[$L^p$空间]
设 $f(x)$ 是 $E\subset\mathbb{R}^n$ 上的可测函数，记
\begin{align*}
\|f\|_p = \left(\int_E |f(x)|^p \mathrm{d}x\right)^{\frac{1}{p}},\quad 0 < p < +\infty.
\end{align*}
并称
\begin{align*}
L^p\left( E \right) =\left\{ f\mid \left\| f \right\| _p<+\infty \right\} 
\end{align*}
为 $\boldsymbol{L}^{\boldsymbol{p}}$ \textbf{空间}.

若存在$M$，使得 $|f(x)| \leqslant M$，a.e., $x\in E$，则称 $f(x)$ 在 $E$ 上\textbf{本性有界}，$M$ 称为 $f(x)$ 的\textbf{本性上界}。
称
\begin{align*}
\left\| f \right\| _{\infty}=\mathrm{inf}\left\{ M\mid \left| f\left( x \right) \right|\leqslant M,\,\mathrm{a}.\mathrm{e}.\,x\in E \right\} 
\end{align*}
为 $f(x)$ 在 $E$ 上的\textbf{本性上确界}。并称
\begin{align*}
L^{\infty}\left( E \right) =\left\{ f\mid f\text{为}E\text{上本性有界函数} \right\} .
\end{align*}
为$\boldsymbol{L}^{\boldsymbol{\infty}}$ \textbf{空间}.
\end{definition}
\begin{remark}
注意，对 $L^p(E)$，我们主要的兴趣在 $p\geqslant 1$ 的情形，下文中若未指明 $p>0$，则一律认为是 $p\geqslant 1$.
\end{remark}

\begin{proposition}\label{proposition:fp极限就是f无穷}
设 $f(x)$ 是 $E\subset\mathbb{R}^n$ 上的可测函数，若 $0 < m(E) < +\infty$，则 $\lim\limits_{p\to\infty} \|f\|_p = \|f\|_\infty$。
\end{proposition}
\begin{proof}
事实上，不妨记 $M = \|f\|_\infty$，则对任一满足 $M' < M$ 的 $M'$，点集
\begin{align*}
A = \{x\in E: |f(x)| > M'\}
\end{align*}
有正测度。由不等式
\begin{align*}
\|f\|_p = \left(\int_E |f(x)|^p \mathrm{d}x\right)^{\frac{1}{p}} \geqslant M'(m(A))^{\frac{1}{p}}
\end{align*}
可知 $\varliminf\limits_{p\to\infty} \|f\|_p \geqslant M'$。令 $M'\to M$，得 $\varliminf\limits_{p\to\infty} \|f\|_p \geqslant M$。

另一方面，我们总有
\begin{align*}
\|f\|_p \leqslant \left(\int_E M^p \mathrm{d}x\right)^{\frac{1}{p}} = M(m(E))^{\frac{1}{p}},
\end{align*}
从而又得
\begin{align*}
\varlimsup\limits_{p\to\infty} \|f\|_p \leqslant M.
\end{align*}
\end{proof}

\begin{theorem}
若 $f,g\in L^p(E)$, $0 < p \leqslant +\infty$, $\alpha,\beta$ 是实数，则
\begin{align*}
\alpha f + \beta g \in L^p(E).
\end{align*}
进而$L^p(E)$构成一个线性空间.
\end{theorem}
\begin{proof}
(i) 当 $0 < p < \infty$ 时，我们有
\begin{align*}
|\alpha f(x) + \beta g(x)|^p \leqslant 2^p(|\alpha|^p \cdot |f(x)|^p + |\beta|^p \cdot |g(x)|^p).
\end{align*}

(ii) 当 $p = +\infty$ 时，我们有
\begin{align*}
|\alpha f(x) + \beta g(x)| \leqslant |\alpha| \cdot \|f\|_\infty + |\beta| \cdot \|g\|_\infty,\quad \text{a.e. } x\in E,
\end{align*}
从而可知 $\|\alpha f + \beta g\|_\infty \leqslant |\alpha| \cdot \|f\|_\infty + |\beta| \cdot \|g\|_\infty$.
\end{proof}

\begin{example}
设 $E=(0,1)$，则
\begin{gather*}
\ln\frac{1}{x} \in L^p(E),\quad \ln\frac{1}{x} \notin L^\infty(E); \\
x^{-\frac{1}{p}} \in L^{p-\alpha}(E)\quad (0 < \alpha < p),\quad x^{-\frac{1}{p}} \notin L^p(E); \\
x^{-\frac{1}{p}}\left(\ln\frac{1}{x}\right)^{-2/p} \notin L^{p+\alpha}(E)\quad (\alpha > 0), \\
x^{-\frac{1}{p}}\left(\ln\frac{1}{x}\right)^{-2/p} \in L^p(E).
\end{gather*}
\end{example}

\begin{definition}
\begin{enumerate}[(1)]
\item 对于 \(1 \leqslant p < \infty\)，\(l^p\) 表示所有满足级数 \(\sum\limits_{n=1}^{\infty} |x_n|^p < \infty\) 的数列 \(x=(x_1,x_2,\cdots)\) 组成的空间，其范数定义为 \(\|x\|_p \triangleq \left( \sum\limits_{n=1}^{\infty} |x_n|^p \right)^{\frac{1}{p}}\)。

\item \(l^{\infty}\) 表示所有有界数列 \(x=(x_1,x_2,\cdots)\) 组成的空间，其上赋予上确界范数 \(\|x\|_{\infty} \triangleq \sup\limits_{n \geqslant 1} |x_n|\)。

\item \(c\) 表示所有收敛数列 \(x=(x_1,x_2,\cdots)\) 组成的空间（即极限 \(\lim\limits_{n \to \infty} x_n\) 存在），范数同样取\(\|x\|_{\infty} \triangleq \sup\limits_{n \geqslant 1} |x_n|\)。

\item \(c_0\) 表示所有收敛到 \(0\) 的数列 \(x=(x_1,x_2,\cdots)\) 组成的空间（即 \(\lim\limits_{n \to \infty} x_n = 0\)），范数同样取\(\|x\|_{\infty} \triangleq \sup\limits_{n \geqslant 1} |x_n|\)。
\end{enumerate}
\end{definition}

\begin{proposition}
我们有
\begin{align*}
c_0 \subset c \subset l^{\infty}.
\end{align*}
在赋予范数\(\|x\|_{\infty} \triangleq \sup\limits_{n \geqslant 1} |x_n|\)的情形下，\(c\) 是 \(l^{\infty}\) 的闭子空间，\(c_0\) 是 \(c\) 的闭子空间。
\end{proposition}

\begin{proposition}
设 \((X, \Sigma, \mu)\) 为测度空间.
\begin{enumerate}[(1)]
\item 如果 \(\mu(X) < \infty\)，则对于 \(0 < q \leqslant p \leqslant +\infty\)，有 \(L^p(X, \mu) \subseteq L^q(X, \mu)\)且$\underset{p\rightarrow \infty}{\lim}\left\| f \right\| _p=\left\| f \right\| _{\infty}.$

\item 对于 \(l^p\) 空间，如果 \(0 < p \leqslant q \leqslant +\infty\)，则有 \(l^p \subseteq l^q\)且$\underset{p\rightarrow \infty}{\lim}\left\| x \right\| _p=\left\| x \right\| _{\infty}.$

\item 如果 \(\mu(X) = \infty\)，则 \(L^p\) 与 \(L^q\) 之间一般不存在包含关系.
\end{enumerate}
\end{proposition}
\begin{proof}
\begin{enumerate}[(1)]
\item $p=q$时 \(L^p(X, \mu) = L^q(X, \mu)\)显然成立.下设$q<p$且\(f \in L^p(X).\)

若 \(p = \infty\)，因为 \(\mu(X) < \infty\)，由 \(L^\infty\) 范数的定义可知，存在常数 \(M = \|f\|_\infty\) 使得 \(|f(x)| \leqslant M\) 几乎处处成立，故
\begin{align*}
\int_X |f|^q \mathrm{d}\mu \leqslant M^q \mu(X) < \infty\Longrightarrow L^\infty \subseteq L^p.
\end{align*}
若 \(p < \infty\)，令 
\begin{align*}
E=\{x\in X:|f(x)|\geqslant 1\},\quad F=\{x\in X:|f(x)|<1\}.
\end{align*}
由于 \(\mu(X) < \infty\)，可知 \(\mu(E) \leqslant \mu(X) < \infty\)，且 \(\mu(F) \leqslant \mu(X) < \infty\)。在 \(E\) 上，因为 \(q < p\)，有 \(|f|^q \leqslant |f|^p\)；在 \(F\) 上，有 \(|f|^q \leqslant 1.\)
因此
\begin{align*}
\int_X |f|^q \mathrm{d}\mu = \int_E |f|^q \mathrm{d}\mu + \int_F |f|^q \mathrm{d}\mu \leqslant \int_E |f|^p \mathrm{d}\mu + \int_F 1 \mathrm{d}\mu \leqslant \int_X |f|^p \mathrm{d}\mu + \mu(X) < \infty.
\end{align*}
所以 \(f \in L^q(X)\)，即 \(L^p \subseteq L^q.\)再由\cref{proposition:fp极限就是f无穷}知$\underset{p\rightarrow \infty}{\lim}\left\| f \right\| _p=\left\| f \right\| _{\infty}.$

\item 设序列 \(x = (x_1, x_2, \cdots) \in l^p\)，即 \(\sum\limits_{n=1}^{\infty} |x_n|^p < \infty.\)
若 \(q = \infty\)，因为 \(|x_n|^p \leqslant \sum\limits_{k=1}^{\infty} |x_k|^p < \infty\)，所以 \(|x_n| \leqslant (\sum |x_k|^p)^{\frac{1}{p}} = \|x\|_p\)，故 \(x \in l^\infty.\)
若 \(q < \infty\)，由于级数 \(\sum\limits_{n=1}^{\infty} |x_n|^p\) 收敛，其通项趋于零，故存在 \(N\)，当 \(n > N\) 时 \(|x_n| < 1\). 此时因为 \(p \leqslant q\)，有 \(|x_n|^q \leqslant |x_n|^p.\)
于是
\begin{align*}
\sum\limits_{n=1}^{\infty} |x_n|^q = \sum\limits_{n=1}^{N} |x_n|^q + \sum\limits_{n=N+1}^{\infty} |x_n|^q \leqslant \sum\limits_{n=1}^{N} |x_n|^q + \sum\limits_{n=N+1}^{\infty} |x_n|^p < \infty.
\end{align*}
所以 \(x \in l^q\)，即 \(l^p \subseteq l^q.\)显然有$\underset{p\rightarrow \infty}{\lim}\left\| x \right\| _p=\left\| x \right\| _{\infty}.$

\item 令 \(X = [1, +\infty)\)，\(f(x) = \frac{1}{x}\). 计算得：
\begin{align*}
\int_{1}^{+\infty} |f(x)|^2 \mathrm{d}x = \int_{1}^{+\infty} \frac{1}{x^2} \mathrm{d}x = 1 < \infty \Longrightarrow f \in L^2([1, +\infty)).
\end{align*}
但是：
\begin{align*}
\int_{1}^{+\infty} |f(x)| \mathrm{d}x = \int_{1}^{+\infty} \frac{1}{x} \mathrm{d}x = \infty \Longrightarrow f \notin L^1([1, +\infty)).
\end{align*}
这说明 \(L^2 \not\subseteq L^1.\)
再定义函数 \(f(x)\)为
\begin{align*}
f(x)\triangleq \begin{cases}
n^2,n\leqslant x<n+\frac{1}{n},\\
0,n+\frac{1}{n^4}\leqslant x<n+1.\\
\end{cases}
\end{align*}
计算得
\begin{align*}
\int_{1}^{+\infty} |f(x)| \mathrm{d}x = \sum\limits_{n=1}^{\infty} n^2 \cdot \frac{1}{n^4} = \sum\limits_{n=1}^{\infty} \frac{1}{n^2} < \infty \Longrightarrow f \in L^1([1, +\infty)).
\end{align*}
但是：
\begin{align*}
\int_{1}^{+\infty} |f(x)|^2 \mathrm{d}x = \sum\limits_{n=1}^{\infty} n^4 \cdot \frac{1}{n^4} = \sum\limits_{n=1}^{\infty} 1 = \infty \Longrightarrow f \notin L^2([1, +\infty)).
\end{align*}
这说明 \(L^1 \not\subseteq L^2.\) 综合上述两例，在测度无限的空间中没有包含关系.
\end{enumerate}
\end{proof}

\begin{definition}[共轭指标]
若 $p,p'>1$, 且 $\frac{1}{p}+\frac{1}{p'}=1$, 则称 $p$ 与 $p'$ 为\textbf{共轭指标(数)}.

若 $p=1$, 则规定共轭指标 $p'=\infty$; 若 $p=\infty$, 则规定共轭指标 $p'=1$.
\end{definition}
\begin{remark}
注意到 $p'=\frac{p}{p-1}$, 可知 $p=2$ 时 $p'=2$.
\end{remark}

\begin{theorem}[Hölder不等式]\label{theorem:实变函数--Hölder不等式}
设 $p$ 与 $p'$ 为共轭指标, 若 $f\in L^p(E), g\in L^{p'}(E)$, 则有
\begin{align}\label{eq::--9ru90RFUJIJR3ruu0-1}
\|fg\|_1 \leqslant \|f\|_p \cdot \|g\|_{p'}, \quad 1\leqslant p\leqslant +\infty,
\end{align}
即
\begin{align*}
\int_E |f(x)g(x)|\mathrm{d}x \leqslant \left( \int_E |f(x)|^p \mathrm{d}x \right)^{\frac{1}{p}} \left( \int_E |g(x)|^{p'} \mathrm{d}x \right)^{\frac{1}{p'}}, \quad 1<p<+\infty,
\end{align*}
以及
\begin{gather*}
\int_E |f(x)g(x)|\mathrm{d}x \leqslant \|g\|_\infty \int_E |f(x)|\mathrm{d}x, \quad p=1;\\
\int_E |f(x)g(x)|\mathrm{d}x \leqslant \|f\|_\infty \int_E |g(x)|\mathrm{d}x, \quad p'=1.
\end{gather*}
特别地,我们有
\begin{enumerate}[(1)]
\item 当 $1<p<+\infty$时，\eqref{eq::--9ru90RFUJIJR3ruu0-1} 中等号成立当且仅当存在常数 $\lambda>0$，使得
\[
|f(x)|^p = \lambda |g(x)|^{p'} \quad \text{a.e. } x\in E,
\]

\item 当 $p=1$时，\eqref{eq::--9ru90RFUJIJR3ruu0-1} 中等号成立当且仅当
\[
|g(x)|=\|g\|_\infty \quad \text{a.e. 在 } \{x\in E: f(x)\neq 0\}\text{ 上}.
\]

\item 当 $p=+\infty$时，\eqref{eq::--9ru90RFUJIJR3ruu0-1} 中等号成立当且仅当
\[
|f(x)|=\|f\|_\infty \quad \text{a.e. 在 } \{x\in E: g(x)\neq 0\}\text{ 上}.
\]
\end{enumerate}
\end{theorem}
\begin{proof}

当 $p$ 或 $p'$ 之一为 $\infty$ 时, \eqref{eq::--9ru90RFUJIJR3ruu0-1}式显然成立.

当 $\|f\|_p=0$ 或 $\|g\|_{p'}=0$ 时, 有 $f(x)g(x)=0$, a.e. $x\in E$, \eqref{eq::--9ru90RFUJIJR3ruu0-1}式也显然成立.

当 $\|f\|_p>0$, $\|g\|_{p'}>0$, 且 $p,p'<+\infty$ 时, 则在公式
\begin{align*}
a^{\frac{1}{p}}b^{\frac{1}{p}'}\leqslant \frac{a}{p}+\frac{b}{p'}, \quad a>0, b>0 \one
\end{align*}
中, 令
\begin{align*}
a=\frac{|f(x)|^p}{\|f\|_p^p}, \quad b=\frac{|g(x)|^{p'}}{\|g\|_{p'}^{p'}},
\end{align*}
可知
\begin{align*}
\frac{|f(x)g(x)|}{\|f\|_p \cdot \|g\|_{p'}} \leqslant \frac{1}{p} \cdot \frac{|f(x)|^p}{\|f\|_p^p} + \frac{1}{p'} \cdot \frac{|g(x)|^{p'}}{\|g\|_{p'}^{p'}}.
\end{align*}
将上式作积分, 即得 $\|fg\|_1\leqslant \|f\|_p \cdot \|g\|_{p'}$.
\end{proof}

\begin{theorem}[Hölder不等式特例:Cauchy-Schwarz不等式]\label{theorem:Hölder不等式特例:Schwarz不等式}
当 $p=p'=2$ 时, 有
\begin{align*}
\int_E |f(x)g(x)|\mathrm{d}x \leqslant \left( \int_E |f(x)|^2 \mathrm{d}x \right)^{\frac{1}{2}} \left( \int_E |g(x)|^2 \mathrm{d}x \right)^{\frac{1}{2}}.
\end{align*}
等号成立当且仅当存在常数 $\lambda>0$，使得
\[
|f(x)|^p = \lambda |g(x)|^{p'} \quad \text{a.e. } x\in E,
\]
\end{theorem}
\begin{remark}
Hölder 不等式对 $\|f\|_p$ 或 $\|g\|_{p'}=+\infty$ 时自然成立.
\end{remark}

\begin{example}
若 $m(E)<+\infty$, 且 $0<p_1<p_2\leqslant+\infty$, 则 $L^{p_2}(E)\subseteq L^{p_1}(E)$, 且有
\begin{align}
\|f\|_{p_1} \leqslant (m(E))^{(\frac{1}{p}_1)-(\frac{1}{p}_2)} \|f\|_{p_2}.
\end{align}
\end{example}

\begin{proof}
不妨设 $p_2<+\infty$. 令 $r=p_2/p_1$, 则 $r>1$. 记 $r'$ 为 $r$ 的共轭指标, 则对 $f\in L^{p_2}(E)$, 由 \eqref{eq::--9ru90RFUJIJR3ruu0-1}式可得
\begin{align*}
&\int_E |f(x)|^{p_1} \mathrm{d}x = \int_E (|f(x)|^{p_1} \cdot 1) \mathrm{d}x \\
&\leqslant \left( \int_E |f(x)|^{p_1 \cdot r} \mathrm{d}x \right)^{1/r} \left( \int_E 1^{r'} \mathrm{d}x \right)^{1/r'} \\
&= (m(E))^{1/r'} \left( \int_E |f(x)|^{p_2} \mathrm{d}x \right)^{1/r},
\end{align*}
从而可知
\begin{align*}
\left( \int_E |f(x)|^{p_1} \mathrm{d}x \right)^{\frac{1}{p}_1} \leqslant (m(E))^{(\frac{1}{p}_1)-(\frac{1}{p}_2)} \left( \int_E |f(x)|^{p_2} \mathrm{d}x \right)^{\frac{1}{p}_2}.
\end{align*}
这就是 \eqref{eq::--9ru90RFUJIJR3ruu0-1}式.
\end{proof}

\begin{example}
若 $f\in L^r(E)\cap L^s(E)$, 且令 $0<r<p<s\leqslant+\infty$ 以及 $0<\lambda<1$, $\frac{1}{p}=\frac{\lambda}{r}+\frac{1-\lambda}{s}$, 则
\begin{align*}
\|f\|_p \leqslant \|f\|_r^\lambda \cdot \|f\|_s^{1-\lambda}.
\end{align*}
\end{example}

\begin{proof}
当 $r<s<+\infty$ 时, 我们有
\begin{align*}
\int_E |f(x)|^p \mathrm{d}x = \int_E |f(x)|^{\lambda p} \cdot |f(x)|^{(1-\lambda)p} \mathrm{d}x 
\leqslant \left( \int_E |f(x)|^r \mathrm{d}x \right)^{\lambda \frac{p}{r}} \left( \int_E |f(x)|^s \mathrm{d}x \right)^{\frac{(1-\lambda)p}{s}}.
\end{align*}
当 $r<s=+\infty$ 时, 因为 $p=r/\lambda$, 所以有
\begin{align*}
\int_E |f(x)|^p \mathrm{d}x \leqslant \|f^{p-r}\|_\infty \int_E |f(x)|^r \mathrm{d}x = \|f\|_\infty^{p\lambda} \cdot \|f\|_r^{p(1-\lambda)}.
\end{align*}
上述结果还说明, 当 $r<p<s\leqslant+\infty$ 时, 有
\begin{align*}
\|f\|_p \leqslant \max\{ \|f\|_r, \|f\|_s \}.
\end{align*}
\end{proof}

\begin{example}
设 $0<r<p<s<+\infty$, $f\in L^p(E)$, 则对任意的 $t>0$, 存在分解 $f(x)=g(x)+h(x)$, 使得
\begin{align*}
\|g\|_r^r \leqslant t^{r-p} \|f\|_p^p, \quad \|h\|_s^s \leqslant t^{s-p} \|f\|_p^p.
\end{align*}
\end{example}

\begin{proof}
对 $t>0$, 作函数
\begin{align*}
g(x) =
\begin{cases}
0, & |f(x)|\leqslant t, \\
f(x), & |f(x)|>t,
\end{cases}
\quad h(x)=f(x)-g(x).
\end{align*}
我们有(注意 $r-p<0$)
\begin{align*}
\|g\|_r^r &= \int_E |g(x)|^r \mathrm{d}x = \int_E |g(x)|^{r-p} |g(x)|^p \mathrm{d}x \\
&\leqslant t^{r-p} \int_E |g(x)|^p \mathrm{d}x \leqslant t^{r-p} \int_E |f(x)|^p \mathrm{d}x,
\end{align*}
即得所证.
类似地可证第二个不等式.
\end{proof}

\begin{theorem}[反向 Hölder 不等式]\label{theorem:反向Hölder不等式}
设 $0<p<1$, $q<0$, $\frac{1}{p}+\frac{1}{q}=1$, 则对 $f\in L^p(E)$, $g\in L^q(E)$, 有
\begin{align*}
\int_E |f(x)g(x)|\mathrm{d}x \geqslant \|f\|_p \cdot \|g\|_q.
\end{align*}
等号成立当且仅当存在常数 $\lambda>0$，使得
\[
|f(x)|^{1-p}\,|g(x)| = \lambda \quad \text{a.e. } x\in E,
\]
这也等价于存在常数 $\lambda>0$，使得
\[
|f(x)| = \lambda \, |g(x)|^{q-1} \quad \text{a.e. } x\in E.
\]
\end{theorem}
\begin{proof}
不妨假定 $fg\in L(E)$, 且令 $\overline{p}=\frac{1}{p}>1$, $\overline{q}=\frac{1}{1-p}>1$, 则 $\frac{1}{\overline{p}}+\frac{1}{\overline{q}}=1$, 且有
\begin{align*}
\|f\|_p = \left( \int_E \frac{|f(x)g(x)|^p}{|g(x)|^p} \mathrm{d}x \right)^{\frac{1}{p}} 
\leqslant \left( \int_E |f(x)g(x)|\mathrm{d}x \right) \left( \int_E \frac{1}{|g(x)|^{\frac{p}{1-p}}} \mathrm{d}x \right)^{\frac{1-p}{p}} 
= \int_E \frac{|f(x)g(x)|}{\|g\|_q} \mathrm{d}x.
\end{align*}
由此即得所证.
\end{proof}

\begin{theorem}[Minkowski不等式]\label{theorem:实变函数--Minkowski不等式}
若 $f,g\in L^p(E)$ ($1\leqslant p\leqslant\infty$), 则
\begin{align}\label{eq::--9ru90RFUJIJR3ruu0-3}
\|f+g\|_p \leqslant \|f\|_p + \|g\|_p.
\end{align}
特别地,我们有
\begin{enumerate}[(1)]
\item 当 $1<p<+\infty$时，\eqref{eq::--9ru90RFUJIJR3ruu0-3} 中等号成立当且仅当存在常数 $\lambda \geqslant 0$且不全为零，使得
\[
f(x)=\lambda g(x)\quad \text{a.e. } x\in E.
\]

\item 当 $p=1$ 时，\eqref{eq::--9ru90RFUJIJR3ruu0-3} 中等号成立当且仅当
\[
f(x)g(x)\geqslant 0\quad \text{a.e. } x\in E.
\]

\item 当 $p=+\infty$ 时，\eqref{eq::--9ru90RFUJIJR3ruu0-3} 中等号成立当且仅当存在正测度集 $A\subset E$，使得在 $A$ 上几乎处处有
\[
|f(x)|=\|f\|_\infty,\qquad |g(x)|=\|g\|_\infty,
\]
且 $f$ 与 $g$ 同号（同向）。若 $\|f\|_\infty=0$ 或 $\|g\|_\infty=0$，则等式自动成立。
\end{enumerate}
\end{theorem}
\begin{proof}

当 $p=1$ 时, \eqref{eq::--9ru90RFUJIJR3ruu0-3}式显然成立; 当 $p=\infty$ 时, 因为 $|f(x)|\leqslant\|f\|_\infty$, a.e. $x\in E$, $|g(x)|\leqslant\|g\|_\infty$, a.e. $x\in E$, 所以有
\begin{align*}
|f(x)+g(x)|\leqslant\|f\|_\infty + \|g\|_\infty, \quad \text{a.e. } x\in E.
\end{align*}
从而可知 $\|f+g\|_\infty\leqslant\|f\|_\infty+\|g\|_\infty$, 即 \eqref{eq::--9ru90RFUJIJR3ruu0-3}式成立.

当 $1<p<+\infty$ 时, 我们有
\begin{align*}
\int_E |f(x)+g(x)|^p \mathrm{d}x
&\leqslant \int_E |f(x)+g(x)|^{p-1} \cdot |f(x)+g(x)|\mathrm{d}x \\
&\leqslant \int_E |f(x)+g(x)|^{p-1} \cdot |f(x)|\mathrm{d}x + \int_E |f(x)+g(x)|^{p-1} \cdot |g(x)|\mathrm{d}x.
\end{align*}
用\hyperref[theorem:实变函数--Hölder不等式]{Hölder不等式}于上式右端第一个积分可得
\begin{align*}
\int_E{\left| f\left( x \right) +g\left( x \right) \right|^{p-1}\cdot \left| f\left( x \right) \right|\mathrm{d}}x\leqslant \left( \int_E{\left| f(x)+g(x) \right|^{\left( p-1 \right) \cdot \frac{p}{p-1}}\mathrm{d}x} \right) ^{\frac{p-1}{p}}\left( \int_E{\left| f\left( x \right) \right|^p\mathrm{d}x} \right) ^{\frac{1}{p}}=\left\| f+g \right\| _{p}^{p-1}\cdot \left\| f \right\| _p;
\end{align*}
同理, 对于前式右端第二个积分也可得
\begin{align*}
\int_E{\left| f\left( x \right) +g\left( x \right) \right|^{p-1}\cdot \left| g\left( x \right) \right|\mathrm{d}}x\leqslant \left( \int_E{\left| f(x)+g(x) \right|^{\left( p-1 \right) \cdot \frac{p}{p-1}}\mathrm{d}x} \right) ^{\frac{p-1}{p}}\left( \int_E{\left| g\left( x \right) \right|^p\mathrm{d}x} \right) ^{\frac{1}{p}}=\left\| f+g \right\| _{p}^{p-1}\cdot \left\| g \right\| _p.
\end{align*}
将上面两式代入前式, 即有
\begin{align*}
\|f+g\|_p^p \leqslant \|f+g\|_p^{p-1} \cdot (\|f\|_p + \|g\|_p).
\end{align*}
不妨设 $\|f+g\|_p\neq0$,于是在上式两端用 $\|f+g\|_p^{p-1}$ 除之, 得
\begin{align*}
\|f+g\|_p \leqslant \|f\|_p + \|g\|_p.
\end{align*}
若 $\|f+g\|_p=0$, 原式无所可证.

最后,由\hyperref[theorem:实变函数--Hölder不等式]{Hölder不等式}等号成立的充要条件容易得到\eqref{eq::--9ru90RFUJIJR3ruu0-3}等号成立的充要条件.
\end{proof}

\begin{theorem}[反向 Minkowski 不等式]\label{theorem:反向Minkowski不等式}
设 $0<p<1$, 则对 $f,g\in L^p(E)$, 有
\begin{align*}
\||f|+|g|\|_p \geqslant \|f\|_p + \|g\|_p.
\end{align*}
上式等号成立当且仅当存在$\lambda>0$,使得
\begin{align*}
|f|=\lambda|g|,\,\,\texorpdfstring{a.e.}\,x\in E.
\end{align*}
\end{theorem}
\begin{proof}
不妨假定 $\||f|+|g|\|_p>0$,记$q=\frac{p}{p-1}$,则由\hyperref[theorem:反向Hölder不等式]{反Hölder不等式}可得
\begin{align*}
\||f|+|g|\|_p^p &= \int_E (|f(x)|+|g(x)|)^{p-1} (|f(x)|+|g(x)|) \mathrm{d}x \\
&\geqslant \left( \int_E (|f(x)|+|g(x)|)^{q(p-1)} \right)^{\frac{1}{q}} (\|f\|_p + \|g\|_p) \\
&= \||f|+|g|\|_p^{\frac{1}{p}} (\|f\|_p + \|g\|_p).
\end{align*}
\end{proof}

\begin{proposition}
\begin{enumerate}[(1)]
\item 设 $f\in L^{p_1}(E)$, $g\in L^{p_2}(E)$ ($0<p_1,p_2<+\infty$), 则 $fg\in L^p(E)$ ($p=\frac{p_1p_2}{p_1+p_2}$).

\item 设 $f(x)$ 是 $\mathbb{R}$ 上的可微函数, 而且 $f'\in L^p(\mathbb{R})$ ($p>1$). 若 $f\in L(\mathbb{R})$, 则 $\lim\limits_{|x|\to+\infty} f(x)=0$.

\item (Hanner 不等式) 设 $f,g\in L^p(E)$ ($1\leqslant p<+\infty$), 则
\begin{enumerate}[(i)]
\item $\|f+g\|_p^p + \|f-g\|_p^p \leqslant (\|f\|_p + \|g\|_p)^p + |\|f\|_p - \|g\|_p|^p$ ($p\geqslant2$);
\item $\|f+g\|_p^p + \|f-g\|_p^p \geqslant (\|f\|_p + \|g\|_p)^p + |\|f\|_p - \|g\|_p|^p$ ($1\leqslant p\leqslant2$);
\item $(\|f+g\|_p + \|f-g\|_p)^p + |\|f+g\|_p - \|f-g\|_p|^p \geqslant 2^p (\|f\|_p^p + \|g\|_p^p)$ ($p\geqslant2$);
\item $(\|f+g\|_p + \|f-g\|_p)^p + |\|f+g\|_p - \|f-g\|_p|^p \leqslant 2^p (\|f\|_p^p + \|g\|_p^p)$ ($1\leqslant p\leqslant2$).
\end{enumerate}

\item (Clarkson 不等式) 设 $1<p,q<+\infty$, $\frac{1}{p}+\frac{1}{q}=1$, $f,g\in L^p(E)$, 则
\begin{enumerate}[(i)]
\item $\left\| \frac{f+g}{2} \right\|_p^p + \left\| \frac{f-g}{2} \right\|_p^p \leqslant \frac{\|f\|_p^p + \|g\|_p^p}{2}$ ($p\geqslant2$);
\item $\left\| \frac{f+g}{2} \right\|_p^p + \left\| \frac{f-g}{2} \right\|_p^p \geqslant \frac{\|f\|_p^p + \|g\|_p^p}{2}$ ($1<p\leqslant2$);
\item $\left\| \frac{f+g}{2} \right\|_p^q + \left\| \frac{f-g}{2} \right\|_p^q \geqslant \left( \frac{\|f\|_p^p + \|g\|_p^p}{2} \right)^{q-1}$ ($p\geqslant2$);
\item $\left\| \frac{f+g}{2} \right\|_p^q + \left\| \frac{f-g}{2} \right\|_p^q \leqslant \left( \frac{\|f\|_p^p + \|g\|_p^p}{2} \right)^{q-1}$ ($1<p\leqslant2$).
\end{enumerate}
\end{enumerate}
\end{proposition}
\begin{proof}
\begin{enumerate}[(1)]
\item 令 $\lambda=\frac{1}{p}_1+\frac{1}{p}_2$, $p=\lambda p_1$, $q=\lambda p_2$, 则 $0<p,q<+\infty$, $\frac{1}{p}+1/q=1$. 由题设 $|f|^{1/\lambda}\in L^p(E)$, $|g|^{1/\lambda}\in L^q(E)$. 从而 $|fg|^{1/\lambda}\in L^1(E)$.
\begin{align*}
\int_E |f(x)g(x)|^p \mathrm{d}x = \int_E |f(x)|^{1/\lambda} |g(x)|^{1/\lambda} \mathrm{d}x 
\leqslant \left( \int_E |f(x)|^{p_1} \mathrm{d}x \right)^{\frac{1}{p}} \left( \int_E |g(x)|^{p_2} \mathrm{d}x \right)^{1/q} < +\infty.
\end{align*}

\item 注意到 $f(x+h)-f(x)=\int_x^{x+h} f'(t)\mathrm{d}t$, 可知
\begin{align*}
|f(x+h)-f(x)| &\leqslant \left| \int_x^{x+h} f'(t)\mathrm{d}t \right| \\
&\leqslant \left| \int_x^{x+h} |f'(t)|\mathrm{d}t \right| \leqslant \left( \int_x^{x+h} |f'(t)|^p \mathrm{d}t \right)^{\frac{1}{p}} \cdot |h|^{1-\frac{1}{p}} \\
&\leqslant \|f'\|_p |h|^{1-\frac{1}{p}} \to 0 \quad (|h|\to0).
\end{align*}
这说明 $f(x)$ 在 $\mathbb{R}$ 上一致连续. 由于 $f\in L(\mathbb{R})$, 故知结论成立.

\item 

\item 
\end{enumerate}
\end{proof}

\begin{example}
设 $f\in L^p([0,1])$ ($p>0$), 则
\begin{align*}
\lim_{p\searrow0} \|f\|_p = \exp\left( \int_0^1 \ln|f(x)|\mathrm{d}x \right).
\end{align*}
\end{example}
\begin{proof}
根据Jensen不等式, 可知
\begin{align*}
\ln\|f\|_p = \frac{1}{p} \ln\left( \int_0^1 |f(x)|^p \mathrm{d}x \right) \geqslant \int_0^1 \ln|f(x)|\mathrm{d}x.
\end{align*}
另一方面, 因为 $\ln x\leqslant x-1$ ($x>0$), 所以有
\begin{align*}
\frac{1}{p} \ln\left( \int_0^1 |f(x)|^p \mathrm{d}x \right) &\leqslant \frac{1}{p} \left( \int_0^1 |f(x)|^p \mathrm{d}x - 1 \right) \\
&= \int_0^1 \frac{|f(x)|^p - 1}{p} \mathrm{d}x.
\end{align*}
令 $E=\{x\in[0,1]: |f(x)|\geqslant1\}$, 我们有: $A_p(x)=[|f(x)|^p-1]/p$ 作为 $p$ 的函数, 在 $E$ 上随 $p\searrow0$ 递减趋于 $\ln|f(x)|$; 在 $[0,1]\setminus E$ 上随 $p\searrow0$ 递增趋于 $\ln|f(x)|$. 由此可得
\begin{align*}
\lim_{p\searrow0} \int_0^1 \frac{|f(x)|^p - 1}{p} \mathrm{d}x = \int_0^1 \ln|f(x)|\mathrm{d}x.
\end{align*}
综合上述结果, 即得所证.
\end{proof}


































\end{document}